\documentclass[11pt]{article}
\usepackage[a4paper,margin=21mm]{geometry}
\usepackage[no-math]{fontspec}
\setmainfont{Latin Modern Roman}
\newfontfamily\CJKfont{Noto Serif CJK SC}
\XeTeXlinebreaklocale "zh"
\XeTeXlinebreakskip=0pt plus 1pt
\usepackage{amsmath,amssymb,amsthm,amscd,graphicx,color}
\usepackage{xurl}
\usepackage[unicode,hidelinks]{hyperref}
\allowdisplaybreaks
\setlength\emergencystretch{3em}
\numberwithin{equation}{section}

\numberwithin{equation}{section}

\newtheorem{Theorem}{Theorem}[section]
\newtheorem{Corollary}[Theorem]{Corollary}
\newtheorem{Lemma}[Theorem]{Lemma}
\newtheorem{Proposition}[Theorem]{Proposition}
 { \theoremstyle{definition}
\newtheorem{Definition}[Theorem]{Definition}
\newtheorem{Note}[Theorem]{Note}
\newtheorem{Example}[Theorem]{Example}
\newtheorem{Remark}[Theorem]{Remark} }

\newcommand{\om}{\omega}

\newcommand{\ti}{\mathrm i}
\newcommand{\ct}{\operatorname{CT}}

\makeatletter\newlabel{ois}{{3}{}{Original source locator}{}{}}\makeatother
\begin{document}
\begin{center}\Large The two Rogers–Ramanujan series-product identities for complex q with absolute value less than one\end{center}
\noindent\textbf{Stable ID:} 990\quad\textbf{Author:} Hjalmar Rosengren\par
\noindent\textbf{Work:} A New (But Very Nearly Old) Proof of the Rogers–Ramanujan Identities\par
\noindent\textbf{Source:} \url{https://sigma-journal.com/2024/059/}\par
\noindent\textbf{Version:} Official SIGMA 20 (2024) article 059, DOI 10.3842/SIGMA.2024.059; journal PDF and native TeX acquired 2026-09-30, exact bytes pinned by SHA256\par
\noindent\textbf{Rights:} CC BY 4.0. Authors retain copyright. \url{https://creativecommons.org/licenses/by/4.0/}. Reuse and typesetting adaptation; original language and mathematical content preserved.\par
\noindent\textbf{Difficulty (prior selection preserved):} {\CJKfont H2: The proof derives quadratic transformations of the generating series by constant-term extraction, rather than assuming the Rogers–Ramanujan evaluations. It then compares a coupled parity recursion for the series and product expressions and establishes coefficient uniqueness. Euler and Jacobi identities are explicit prerequisites, while the nontrivial transformation and matching recursion remain in the author proof; the two companion identities form one unit.}\par
\medskip\noindent\textbf{Editorial boundary note (not author text):} The two original Rogers–Ramanujan identities are one companion outcome, not two counts. Complete original statements, all product/series/theta/Pochhammer definitions and full current constant-term quadratic transformations, parity recursions, product-side theta comparison and coefficient uniqueness argument are retained. The entire separate author presentation of Rogers quadratic-transform proof is also included, preserving all steps and cases. Euler, Jacobi and the displayed Ramanujan summation form are explicitly stated formal standard prerequisites with exact author citations. The source gives one product-side even-part reduction and says the other parts follow similarly, while displaying both general even/odd theta formulas and the alternative direct theta split; this is literal repeated classical theta algebra, not an editor-supplied derivation. The exact coupled coefficient recursions are written, not inferred. Independent Section3 variants are excluded. Literal source signs, subscript conventions and printed formulas remain unchanged; no reproof, mathematical refereeing or formula repair.\par
\noindent\textbf{External original reference ois:} Original additional-identity section, source272–350, PDF6–8. Only a discussion of further unselected byproducts/variants; none is counted separately.\par
\medskip\hrule\medskip\noindent\textbf{Author text begins below.}\par

% Original source lines 67--271
\section{Introduction}
The fascinating story of the Rogers--Ramanujan identities
\begin{subequations}\label{rr}
\begin{align}\sum_{k=0}^\infty\frac{q^{k^2}}{(1-q)\bigl(1-q^2\bigr)\dotsm(1-q^k)}&=\prod_{j=0}^\infty\frac 1{\bigl(1-q^{5j+1}\bigr)\bigl(1-q^{5j+4}\bigr)},\label{rra}\\
\sum_{k=0}^\infty\frac{q^{k(k+1)}}{(1-q)(1-q^2)\dotsm\bigl(1-q^k\bigr)}&=\prod_{j=0}^\infty\frac 1{\bigl(1-q^{5j+2}\bigr)\bigl(1-q^{5j+3}\bigr)}\label{rrb}
\end{align}
\end{subequations}
has been told many times. They were first proved by Rogers in a brilliant but forgotten paper from 1894 and rediscovered by Ramanujan some time before 1913.
Hardy wrote about the seven proofs known in the 1930s: ``None of them is both simple and straight-forward, and probably it would be unreasonable to hope for a proof which is'' \cite{h}.
Several dozen more proofs have appeared since then, using diverse methods.
We refer to \cite{a89} for a survey of the proofs available in the 1980s and \cite{akp,bp,c,f,gis,p,st,t} for examples of more recent work.
The purpose of the present note is to give a new proof that, although far from the shortest possible, we find relatively ``simple and straight-forward''.
Almost unbelievably, this proof was very nearly found by Rogers. He knew all its main ingredients and could easily have given it by adding
one or two sentences to \cite{r2}.
Among the many gems buried in Rogers' paper are the identities \cite[Section~6, formulas~(5)--(8)]{r2}
\begin{subequations}\label{ghr}
\begin{alignat}{3}%\label{ghra}
&G(q)+G(-q)=2F(q)G\bigl(q^{16}\bigr),\qquad&&
G(q)-G(-q)=2qF(q)H\bigl(-q^4\bigr),&\\
&H(q)+H(-q)=2F(q)G\bigl(-q^4\bigr),\qquad&&
H(q)-H(-q)=2q^3F(q)H\bigl(q^{16}\bigr),&
\end{alignat}
\end{subequations}
where $G(q)$ and $H(q)$ denote the sums in \eqref{rra} and \eqref{rrb}, respectively, and
\begin{equation*}
F(q)=\prod_{j=0}^\infty\frac{1-q^{8j+8}}{1-q^{2j+2}}.
\end{equation*}
 Below, we will give both a version of Rogers' proof of \eqref{ghr} and another proof based on work of Andrews \cite{a86}.
After deriving \eqref{ghr}, Rogers writes somewhat dismissively
\begin{quote}These four identities \dots\ are sufficiently remarkable in themselves to call for mention at this point, although they may all be derived from the $\theta$-function values of the series $\phi(q)$, $\psi(q)$ [our $G(q)$ and $H(q)$] obtained in the last section. [\dots] It is not, however, in these identities that the special interest in the series $\phi(q)$ and $\psi(q)$ lies.
\end{quote}
At this point Rogers has already proved the Rogers--Ramanujan identities. He realizes that
 \eqref{ghr} are quite easy to prove for the product sides of \eqref{rr} and hence perhaps less interesting than they may seem.

Our observation is simply that the identities \eqref{ghr} uniquely determine $G$ and $H$ up to normalization. In fact, if they are written as
\begin{subequations}\label{abf}
\begin{gather}
G(q)=F(q)\bigl(G\bigl(q^{16}\bigr)+qH\bigl(-q^4\bigr)\bigr), \\
H(q)=F(q)\bigl(G\bigl(-q^4\bigr)+q^3H\bigl(q^{16}\bigr)\bigr),
\end{gather}
\end{subequations}
 they can be used to recursively compute $G$ and $H$ as power series, starting from the initial values $G(0)=H(0)=1$. Since Rogers has proved that \eqref{ghr} hold both for the sum sides and the product sides of \eqref{rr}, this gives an independent proof of the Rogers--Ramanujan identities, which is
considerably easier than the one given earlier in the same paper.

We will present the resulting proof in some detail, but we stress that apart from the preceding paragraph everything can be found in
 \cite{a86,r2}.
I am very much indebted to George Andrews for pointing this out to me.

In Section~\ref{ois}, we explain how some further known Rogers--Ramanujan-type identities either appear as byproducts of the proof or can be obtained by slight variations.
We only consider identities that are very close to \eqref{rr}, in that the product sides are essentially the same. It would be interesting to investigate whether similar proofs can be given for other results. As a first step,
in \cite{ro} we show that the method can be used to prove
Rogers' and Slater's mod $7$ identities.

\section{A proof of the Rogers--Ramanujan identities}\label{ps}

We will use the standard notation
\begin{equation*}
(a;q)_k=\prod_{j=0}^{k-1}\bigl(1-aq^j\bigr),\qquad k\in\mathbb Z_{\geq 0}\cup\{\infty\}, \end{equation*}
which is extended to negative indices by $(a;q)_{-k}=1/\bigl(aq^{-k};q\bigr)_k$. More generally,
\begin{equation*}(a_1,\dots,a_m;q)_k=(a_1;q)_k\dotsm (a_m;q)_k. \end{equation*}
We will write
\begin{equation*}R(t;q)=\sum_{k=0}^\infty\frac{q^{k^2}t^k}{(q;q)_k} \end{equation*}
and $G(q)=R(1;q)$, $H(q)=R(q;q)$.
The Rogers--Ramanujan identities \eqref{rr} then take the compact form
\[%\label{rrgh}
G(q)=\frac 1{\bigl(q,q^4;q^5\bigr)_\infty},\qquad H(q)=\frac 1{\bigl(q^2,q^3;q^5\bigr)_\infty}.
\]
Finally, we introduce the multiplicative theta function $\theta(z;q)=(q,z,q/z;q)_\infty$.

The only results that we will need are Euler's identity
\begin{equation}\label{esa}\frac 1{(z;q)_\infty}=\sum_{k=0}^\infty\frac{z^k}{(q;q)_k},\qquad |z|<1,\end{equation}
 Jacobi's identity
\begin{gather}\theta(z;q)=\sum_{k=-\infty}^\infty(-1)^kq^{k(k-1)/2}z^k \label{jti}
\end{gather}
and
\begin{gather}
\frac{\theta(1/z;q)}{(tz;q)_\infty}=(t;q)_\infty\sum_{k=-\infty}^{\infty}\frac{(-1)^kq^{k(k-1)/2}}{(t;q)_kz^k}, \qquad 0<|z|<\frac 1{|t|}.\label{rps}
\end{gather}
To some readers these identities may look just as complicated as \eqref{rr}, but they are much easier to prove. For instance, they can all be obtained as special cases of Ramanujan's ${}_1\psi_1$-summation, a simple proof of which was given by Ismail \cite{i}, see also \cite[Appendix~C]{a86}.


Following Andrews \cite{a86}, we will represent the sums in \eqref{rr} as constant terms. In general, if~${f(z)=\sum_{k=0}^\infty a_k z^k}$
 (all series in this paper can be considered as formal or analytic according to taste)
and we let~$\ct$ denote the constant term of a Laurent series in $z$, then by \eqref{jti}, we~have%
\begin{align}
\ct \theta(1/z;Q)f\bigl(xz^N\bigr)&=\ct \sum_{l=-\infty}^\infty (-1)^lQ^{l(l-1)/2}z^{-l}\sum_{k=0}^\infty a_k x^kz^{Nk}\nonumber\\
&=\sum_{k=0}^\infty
(-1)^{Nk}Q^{Nk(Nk-1)/2}a_k x^k.\label{fexp}
\end{align}
We apply this with $f(z)$ as Euler's series \eqref{esa}. In order to get a factor $q^{k^2}$, we must choose $Q=q^{2/N^2}$. We will only need the cases $N=1$ and $N=2$, which give the constant term representations
\begin{equation*}R(t;q)
=\ct\frac{\theta\bigl(1/z;q^{2}\bigr)}{ (-qtz;q)_\infty}
=\ct\frac{\theta\bigl(1/z;q^{1/2}\bigr)}{ \bigl(q^{1/2}tz^2;q\bigr)_\infty}.
\end{equation*}

Next, we write
 \begin{gather*}
 \frac{\theta\bigl(1/z;q^{2}\bigr)}{ (-qtz;q)_\infty}= \frac{\theta\bigl(1/z;q^{2}\bigr)}{\bigl(-qtz;q^2\bigr)_\infty}\cdot\frac1{\bigl(-q^2tz;q^2\bigr)_\infty}=\frac{\theta(1/z;q^{2})_\infty}{\bigl(-q^2tz;q^2\bigr)_\infty}\cdot\frac1{\bigl(-qtz;q^2\bigr)_\infty},\\
\frac{\theta\bigl(1/z;q^{1/2}\bigr)}{ \bigl(q^{1/2}tz^2;q\bigr)_\infty}= \frac{\theta\bigl(1/z;q^{1/2}\bigr)}{ \bigl(q^{1/4}t^{1/2}z;q^{1/2}\bigr)_\infty}\cdot \frac 1{ \bigl(- q^{1/4}t^{1/2}z;q^{1/2}\bigr)_\infty},
\end{gather*}
and expand the resulting factors using \eqref{esa} and \eqref{rps}. This gives the quadratic transformation formulas
\begin{subequations}\label{rtf}
\begin{align}
R(t;q)&=\bigl(-qt;q^2\bigr)_\infty\sum_{k=0}^\infty\frac{q^{k(k+1)}t^k}{\bigl(q^2,-qt;q^2\bigr)_k}\\
&=\bigl(-q^2t;q^2\bigr)_\infty\sum_{k=0}^\infty\frac{q^{k^2}t^k}{\bigl(q^2,-q^2t;q^2\bigr)_k}\label{rtfb}\\
&=\bigl(q^{1/4}t^{1/2};q^{1/2}\bigr)_\infty\sum_{k=0}^\infty\frac{q^{k^2/4}t^{k/2}}{\bigl(q^{1/2},q^{1/4}t^{1/2};q^{1/2}\bigr)_k}.\label{rtfc}
\end{align}
\end{subequations}
These identities were found by Rogers \cite[Section~6, formulas~(2)--(4)]{r2}. Except for his way of writing, the original proof is not difficult; we will review it in Section~\ref{rpss}.
 Watson independently proved \eqref{rtfb} and \eqref{rtfc} \cite{w}.

For $t=1$ and $t=q$, the equations \eqref{rtf} simplify to
\begin{subequations}\label{rm20}
\begin{align}
G(q)&=\bigl(-q;q^2\bigr)_\infty\sum_{k=0}^\infty\frac{q^{k(k+1)}}{(-q;-q)_{2k}}\\
& =\bigl(-q^2;q^2\bigr)_\infty\sum_{k=0}^\infty\frac{q^{k^2}}{\bigl(q^4;q^4\bigr)_k} \\
&=\bigl(q^{1/4};q^{1/2}\bigr)_\infty\sum_{k=0}^\infty\frac{q^{{k^2}/4}}{\bigl(q^{1/4};q^{1/4}\bigr)_{2k}},\\
H(q)&=\bigl(-q^2;q^2\bigr)_\infty\sum_{k=0}^\infty\frac{q^{k(k+2)}}{\bigl(q^4;q^4\bigr)_k}\\
& =\bigl(-q;q^2\bigr)_\infty\sum_{k=0}^\infty\frac{q^{k(k+1)}}{(-q;-q)_{2k+1}} \\
&=\bigl(q^{1/4};q^{1/2}\bigr)_\infty\sum_{k=0}^\infty\frac{q^{{k(k+2)}/4}}{\bigl(q^{1/4};q^{1/4}\bigr)_{2k+1}}.
\end{align}
 \end{subequations}

Let us introduce the auxiliary series
\begin{equation}\label{ab}
A(q)=\sum_{k=0}^\infty\frac{q^{k^2}}{\bigl(q^4;q^4\bigr)_k},\qquad B(q)=\sum_{k=0}^\infty\frac{q^{k(k+2)}}{\bigl(q^4;q^4\bigr)_k}. \end{equation}
Then, \eqref{rm20} can be written
\begin{subequations}\label{ghij}
\begin{gather}
 G\bigl(-q^4\bigr)=\bigl(q^4;q^8\bigr)_\infty\frac{B(q)+B(-q)}{2},\label{ga}\\
G(q)=\bigl(-q^2;q^2\bigr)_\infty A(q),\\
 G\bigl(q^{16}\bigr)=\bigl(q^4;q^8\bigr)_\infty\frac{A(q)+A(-q)}{2}, \\
H(q)=\bigl(-q^2;q^2\bigr)_\infty B(q),\label{hb}\\
 H\bigl(-q^4\bigr)=\bigl(q^4;q^8\bigr)_\infty\frac{A(q)-A(-q)}{2q},\\
H\bigl(q^{16}\bigr)=\bigl(q^4;q^8\bigr)_\infty\frac{B(q)-B(-q)}{2q^3}.
\end{gather}
\end{subequations}
We may use \eqref{ga} and \eqref{hb} to eliminate
$A$ and $B$ from the remaining equations. Writing also
\begin{equation*}
F(q)=\frac{\bigl(-q^2;q^2\bigr)_\infty}{\bigl(q^4;q^8\bigr)_\infty}=\frac{\bigl(q^8;q^8\bigr)_\infty}{\bigl(q^2;q^2\bigr)_\infty},
\end{equation*}
we arrive at \eqref{ghr} or, equivalently, \eqref{abf}.

As was noted in the introduction, the identities \eqref{abf}
can be considered as recursions for computing the Taylor series of $G$ and $H$, starting from $G(0)=H(0)=1$; see \eqref{rec} below. Hence, if we can verify that
 \eqref{ghr} also holds for the product sides
\begin{equation*}
\tilde G(q)=\frac 1{\bigl(q,q^4;q^5\bigr)_\infty},\qquad
\tilde H(q)=\frac 1{\bigl(q^2,q^3;q^5\bigr)_\infty},\end{equation*}
then the Rogers--Ramanujan identities $G=\tilde G$, $H=\tilde H$ follow.
As Rogers observed, this reduces to classical theta function identities.
We first write
\begin{align*}
\tilde G(q)&=\frac 1{\bigl(q,q^4,q^6,q^9;q^{10}\bigr)_\infty}=\frac{\bigl(-q,-q^9;q^{10}\bigr)_\infty}{\bigl(q^4,q^6;q^{10}\bigr)_\infty\bigl(q^2,q^{18};q^{20}\bigr)_\infty}=\frac{\bigl(q^8,q^{12};q^{20}\bigr)_\infty \theta\bigl(-q;q^{10}\bigr)}{(q^2;q^{2})_\infty}
 \end{align*}
 and similarly
 \begin{equation*}
 \tilde H(q)=\frac{\bigl(q^4,q^{16};q^{20}\bigr)_\infty \theta\bigl(-q^3;q^{10}\bigr)}{\bigl(q^2;q^{2}\bigr)_\infty}. \end{equation*}
Thus, we are reduced to the problem of computing the even and odd part of a theta function.
In general, Jacobi's triple product identity \eqref{jti} implies the classical identities
\begin{align*}
&\frac{\theta(z;q)+\theta(-z;q)}2=\sum_{k=-\infty}^\infty q^{k(2k-1)}z^{2k}=\theta\bigl(-q z^2;q^4\bigr), \\
&\frac{\theta(z;q)-\theta(-z;q)}2=-\sum_{k=-\infty}^\infty q^{k(2k+1)}z^{2k+1}=-z\theta\bigl(-q^3z^2;q^4\bigr).
\end{align*}
It follows that, for instance,
\begin{equation*}\frac{\tilde G(q)+\tilde G(-q)}2=\frac{\bigl(q^8,q^{12};q^{20}\bigr)_\infty\bigl(-q^{12},-q^{28},q^{40};q^{40}\bigr)_\infty}{\bigl(q^2;q^{2}\bigr)_\infty}.\end{equation*}
Using that
\begin{equation*}
\bigl(-q^{12},-q^{28};q^{40}\bigr)_\infty =\frac {\bigl(q^{24},q^{56};q^{80}\bigr)_\infty}{\bigl(q^{12},q^{28};q^{40}\bigr)_\infty}
=\frac{\bigl(q^{16},q^{24};q^{40}\bigr)_\infty}{\bigl(q^{12},q^{28};q^{40}\bigr)_\infty\bigl(q^{16},q^{64};q^{80}\bigr)},
\end{equation*}
the expression above is easily reduced to
\begin{equation*}\frac{\tilde G(q)+\tilde G(-q)}2 =\frac{\bigl(q^8;q^8\bigr)_\infty}{\bigl(q^2;q^2\bigr)_\infty\bigl(q^{16},q^{64};q^{80}\bigr)_\infty}=\frac{\bigl(q^8;q^8\bigr)_\infty}{\bigl(q^2;q^2\bigr)_\infty} \tilde G\bigl(q^{16}\bigr).
\end{equation*}
The remaining parts of \eqref{ghr}, with $G$ and $H$ replaced by $\tilde G$ and $\tilde H$, follow similarly.
Alternatively, the identities \eqref{abf} can be obtained directly as special cases of
\begin{equation*}
\theta(z;q)=\theta\bigl(-qz^2;q^4\bigr)-z \theta\bigl(-q^3z^2;q^4\bigr).
\end{equation*}

Finally, we make some remarks on the
 recursions for the Taylor coefficients of $G$ and $H$ mentioned above.
Substituting
\begin{equation*}F(q)=\sum_{k=0}^\infty f_kq^{2k},\qquad G(q)=\sum_{k=0}^\infty g_kq^k,\qquad H(q)=\sum_{k=0}^\infty h_kq^k\end{equation*}
into \eqref{abf} gives
\begin{subequations}\label{rec}
\begin{alignat}{3}%\label{reca}
& g_{2k}=\sum_{j=0}^{\lfloor k/8\rfloor}f_{k-8j}g_j,\qquad&&
g_{2k+1}=\sum_{j=0}^{\lfloor k/2\rfloor}(-1)^jf_{k-2j}h_j,&\\
& h_{2k}=\sum_{j=0}^{\lfloor k/2\rfloor}(-1)^jf_{k-2j}g_j,\qquad&&
h_{2k+1}=\sum_{j=0}^{\lfloor (k-1)/8\rfloor}f_{k-8j-1}g_j.&
\end{alignat}
\end{subequations}
It is interesting to note that all the coefficients involved here have partition-theoretic interpretations. This is straight-forward for $f_n$, and very well-known for $g_n$ and $h_n$.
This leads to the following result, which can be viewed as a combinatorial reformulation of~\eqref{abf}.

\begin{Corollary}
Let $f_n$ be the number of partitions of $n$ into parts not divisible by $4$.
Let $g_n$ be the number of partitions of $n$ into parts congruent to $\pm 1$ $\operatorname{mod}$ $5$ or, equivalently, the number of partitions
of $n$ without repeated or consecutive parts. Let $h_n$ be
the number of partitions of $n$ into parts congruent to $\pm 2$ $\operatorname{mod}$ $5$ or, equivalently, the number of partitions
of $n$ without repeated or consecutive parts and not containing $1$ as a part. Then, the recursions \eqref{rec} hold.
\end{Corollary}


\setcounter{section}{3}
% Original source lines 351--412
\section{Rogers' proof of the quadratic transformations}\label{rpss}

Since Rogers' papers are not easy to read, we will provide a slightly modified
version of his proof of the quadratic transformations \eqref{rtf}.
Although the notation was not used by Rogers, we find it instructive to introduce the $q$-exponential functions
\begin{align*}
e_q(x)&=\frac 1{(x;q)_\infty}=\sum_{k=0}^\infty\frac{x^k}{(q;q)_k},\qquad
E_q(x)=(-x;q)_\infty=\sum_{k=0}^\infty\frac{q^{k(k-1)/2}x^k}{(q;q)_k}.
\end{align*}
Roger defines the $q$-derivative as
\begin{equation*}(\delta_x f)(x)=\frac{f(x)-f(qx)}{x}. \end{equation*}
It is easy to see that $\delta_x e_q(tx)=t e_q(tx)$,
 and hence $f(\delta_x)e_q(tx)=f(t)e_q(tx)$ for any formal power series $f$. In particular,
\begin{equation}\label{dxe}
e_q(y\delta_x)e_q(tx)=e_q(tx)e_q(ty).\end{equation}
Identifying the coefficient of $t^k$ in \eqref{dxe} gives
\begin{equation}\label{dxh}e_q(y\delta_x)\frac{x^k}{(q;q)_k}=H_k(x,y),\end{equation}
where $H_k$ are defined by the generating function
\begin{equation}\label{eeh}e_q(tx)e_q(ty)= \frac 1{(tx,ty;q)_\infty}=\sum_{k=0}^\infty H_k(x,y)t^k.\end{equation}
In \cite{r1}, Rogers proves \eqref{dxh} in a more direct way and uses it to derive \eqref{dxe}.

The second $q$-exponential function satisfies $\delta_x E_q(tx)=t E_q(qtx)$,
which gives $\delta_x^k E_q(tx)=q^{k(k-1)/2}t^kE_q\bigl(q^ktx\bigr)$ and
\begin{equation*} e_q(y\delta_x)E_q(tx)=\sum_{k=0}^\infty\frac{q^{k(k-1)/2}y^kt^k}{(q;q)_k}\,E_q(q^ktx)
=(-tx;q)_\infty\sum_{k=0}^\infty\frac{q^{k(k-1)/2}y^kt^k}{(-tx,q;q)_k}.\end{equation*}
On the other hand, \eqref{dxh} gives
\begin{equation*}e_q(y\delta_x)E_q(tx)=\sum_{k=0}^\infty t^kq^{k(k-1)/2} H_k(x,y). \end{equation*}
Specializing $t=q$, we obtain \cite[Section~6, equation~(1)]{r2} in the form
\begin{equation}\label{rhi}\sum_{k=0}^\infty q^{k(k+1)/2} H_k(x,y)=(-qx;q)_\infty\sum_{k=0}^\infty\frac{q^{k(k+1)/2}y^k}{(-qx,q;q)_k}.
\end{equation}
By symmetry, one can interchange $x$ and $y$ on the right-hand side.


Roger now observes that the case $y=xq^{1/2}$ of \eqref{eeh} is
\begin{equation*}
 \sum_{k=0}^\infty H_k\bigl(x,xq^{1/2}\bigr)t^k=\frac 1{\bigl(tx;q^{1/2}\bigr)_\infty},\end{equation*}
which implies
\begin{equation*}
H_k\bigl(x,xq^{1/2}\bigr)=\frac{ x^k}{\bigl(q^{1/2};q^{1/2}\bigr)_k}.\end{equation*}
Using this in \eqref{rhi} and replacing $q$ by $q^2$ gives
\begin{subequations}\label{qt}
\begin{gather}
\sum_{k=0}^\infty \frac {q^{k(k+1)}x^k}{(q;q)_k}=\bigl(-q^2x;q^2\bigr)_\infty\sum_{k=0}^\infty\frac{q^{k(k+2)}x^k}{\bigl(-q^{2}x,q^2;q^2\bigr)_k}.
\end{gather}
Interchanging the role of $x$ and $y$ gives the alternative expression
\begin{gather}
\sum_{k=0}^\infty \frac {q^{k(k+1)}x^k}{(q;q)_k}=\bigl(-q^3x;q^2\bigr)_\infty\sum_{k=0}^\infty\frac{q^{k(k+1)}x^k}{\bigl(-q^{3}x,q^2;q^2\bigr)_k}.
\end{gather}
Finally, the case $y=-x$ of \eqref{eeh} is
\begin{equation*}
\sum_{k=0}^\infty H_k(x,-x)t^k=\frac 1{\bigl(t^2x^2;q^{2}\bigr)_\infty},
\end{equation*}
which gives
\begin{equation*}
H_{2k}(x,-x)=\frac{x^{2k}}{\bigl(q^2;q^2\bigr)_k},\qquad H_{2k+1}(x,x)=0.
 \end{equation*}
Using this in \eqref{rhi} gives
\begin{equation}
\sum_{k=0}^\infty \frac{q^{2k(2k+1)}x^{2k}}{\bigl(q^2;q^2\bigr)_k}=(-qx;q)_\infty\sum_{k=0}^\infty\frac{(-1)^kq^{k(k+1)/2}x^k}{(-qx,q;q)_k}.
\end{equation}
\end{subequations}
The identities \eqref{qt} are indeed equivalent to \eqref{rtf}.
\begin{thebibliography}{99}
\bibitem[1]{a86}
Andrews G.E., {$q$}-series: their development and application in analysis,
 number theory, combinatorics, physics, and computer algebra, \textit{CBMS
 Reg. Conf. Ser. Math.}, Vol.~66, \href{https://doi.org/10.1090/cbms/066}{American Mathematical Society}, Providence,
 RI, 1986.

\bibitem[2]{a89}
Andrews G.E., On the proofs of the {R}ogers--{R}amanujan identities, in
 {$q$}-{S}eries and {P}artitions ({M}inneapolis, {MN}, 1988), \textit{IMA Vol.
 Math. Appl.}, Vol.~18, \href{https://doi.org/10.1007/978-1-4684-0637-5_1}{Springer}, New York, 1989, 1--14.

\bibitem[3]{akp}
Andrews G.E., Knopfmacher A., Paule P., An infinite family of {E}ngel
 expansions of {R}ogers--{R}amanujan type, \href{https://doi.org/10.1006/aama.2000.0686}{\textit{Adv. in Appl. Math.}}
 \textbf{25} (2000), 2--11.

\bibitem[5]{bp}
Boulet C., Pak I., A combinatorial proof of the {R}ogers--{R}amanujan and
 {S}chur identities, \href{https://doi.org/10.1016/j.jcta.2005.09.007}{\textit{J.~Combin. Theory Ser.~A}} \textbf{113} (2006),
 1019--1030, \href{https://arxiv.org/abs/math.CO/0411072}{arXiv:math.CO/0411072}.

\bibitem[6]{c}
Corteel S., Rogers--{R}amanujan identities and the
 {R}obinson--{S}chensted--{K}nuth correspondence, \href{https://doi.org/10.1090/proc/13373}{\textit{Proc. Amer. Math.
 Soc.}} \textbf{145} (2017), 2011--2022.

\bibitem[7]{f}
Fulman J., A probabilistic proof of the {R}ogers--{R}amanujan identities,
 \href{https://doi.org/10.1017/S0024609301008207}{\textit{Bull. London Math. Soc.}} \textbf{33} (2001), 397--407,
 \href{https://arxiv.org/abs/math.CO/0001078}{arXiv:math.CO/0001078}.

\bibitem[8]{gis}
Garrett K., Ismail M.E.H., Stanton D., Variants of the {R}ogers--{R}amanujan
 identities, \href{https://doi.org/10.1006/aama.1999.0658}{\textit{Adv. in Appl. Math.}} \textbf{23} (1999), 274--299.

\bibitem[10]{h}
Hardy G.H., Ramanujan. {T}welve lectures on subjects suggested by his life and
 work, Cambridge University Press, Cambridge, 1940.

\bibitem[11]{i}
Ismail M.E.H., A simple proof of {R}amanujan's {$_{1}\psi_{1}$} sum,
 \href{https://doi.org/10.2307/2041093}{\textit{Proc. Amer. Math. Soc.}} \textbf{63} (1977), 185--186.

\bibitem[13]{p}
Paule P., Short and easy computer proofs of the {R}ogers--{R}amanujan
 identities and of identities of similar type, \href{https://doi.org/10.37236/1190}{\textit{Electron.~J. Combin.}}
 \textbf{1} (1994), R10, 9~pages.

\bibitem[14]{r1}
Rogers L.J., On the expansion of some infinite products, \href{https://doi.org/10.1112/plms/s1-24.1.337}{\textit{Proc. Lond.
 Math. Soc.}} \textbf{24} (1893), 337--352.

\bibitem[15]{r2}
Rogers L.J., Second memoir on the expansion of certain infinite products,
 \href{https://doi.org/10.1112/plms/s1-25.1.318}{\textit{Proc. Lond. Math. Soc.}} \textbf{25} (1893), 318--343.

\bibitem[17]{ro}
Rosengren H., New proofs of the septic {R}ogers--{R}amanujan identities,
 \href{https://arxiv.org/abs/2302.02312}{arXiv:2302.02312}.

\bibitem[21]{st}
Stembridge J.R., {H}all--{L}ittlewood functions, plane partitions, and the
 {R}ogers--{R}amanujan identities, \href{https://doi.org/10.2307/2001250}{\textit{Trans. Amer. Math. Soc.}}
 \textbf{319} (1990), 469--498.

\bibitem[22]{t}
Tsuchioka S., A proof of the second {R}ogers--{R}amanujan identity via
 {K}leshchev multipartitions, \href{https://doi.org/10.3792/pjaa.99.005}{\textit{Proc. Japan Acad. Ser.~A Math. Sci.}}
 \textbf{99} (2023), 23--26, \href{https://arxiv.org/abs/2205.06038}{arXiv:2205.06038}.

\bibitem[23]{w}
Watson G.N., A note on Lerch's functions, \href{https://doi.org/10.1093/qmath/os-8.1.43}{\textit{Q.~J.~Math.}} \textbf{8}
 (1937), 43--47.

\end{thebibliography}
\end{document}
